% H9-T22: wild positive-genus covers and generic hyperelliptic covers in char2.
% Run JAC1964-20261121; /root/h9_height; GPT-6 Astra Pro.

\subsection*{H9-T22. Wild simple-pole covers and every generic hyperelliptic genus in characteristic two}

\begin{theorem}[H9-T22]
Let $k$ be algebraically closed of characteristic $p>0$,
and let $A/k$ and its Lefschetz pencil be as in H9-T21.
Choose the coordinate so that both $0$ and $\infty$ are
smooth pencil values. Choose $b\ge0$, distinct nonzero
$b_1,\ldots,b_b\in k-\Sigma$, and $c_1,\ldots,c_b\in k^\times$.
Set
\[
 a(t)=t+\sum_{j=1}^{b}\frac{c_j}{t-b_j},\qquad
 F=\overline{k(s)},\qquad K=F(t),
 \qquad L=F(t,w),\quad w^p-w=a(t)+\frac{s}{t},
                                                               \tag{H9.117}
\]
where $s$ is transcendental over $k$. Then $L$ is the
function field of a smooth projective curve of genus
$(p-1)(b+1)$, and $L/K$ has degree $p$ and is unramified
above $\Sigma$. On the entire homologically trivial
rational codimension-$m$ cycle group,
\[
 h_L(z,w')=\frac1p h_K(\operatorname{cor}z,
                              \operatorname{cor}w'),\qquad
 \operatorname{rad}h_L
       =\operatorname{cor}^{-1}(\operatorname{rad}h_K).
                                                               \tag{H9.118}
\]
All these cycles are admissible. The same conclusion
holds at the geometric generic point of the full
simple-pole and nonzero-residue parameter stratum
described below.

In characteristic two, for every genus $g\ge1$, the
conclusion with factor $1/2$ holds for the geometric
generic hyperelliptic degree-two cover of $\mathbf P^1$
unramified above $\Sigma$, in the full double-cover
parameter space. Consequently, together with H9-T20,
geometric generic hyperelliptic covers of every genus
in every characteristic satisfy the exact norm-height
formula. Under 032, all of these full homologically
trivial height groups are positive semidefinite when
$\dim A=4$.
\end{theorem}

\begin{proof}
\paragraph{The wild curve family is explicitly smooth.}
Let $S=\mathbf G_{m,s}$. Normalize
$\mathbf P^1_t\times S$ in the Artin--Schreier extension
(H9.117). Its poles are the $b+2$ disjoint sections
$0,b_1,\ldots,b_b,\infty$, all simple with nonzero
residues, and all outside $\Sigma$. Away from these
sections the equation is finite \'etale because its
derivative in $w$ is $-1$.

Here is a uniform local normalization at every pole.
With base uniformizer $z$, write the right-hand side as
$c/z+d(z,s)$, where $c$ is a unit and $d$ is regular.
Put $v=1/w$. The closure near the ramified point has
equation
\[
       z\bigl(1-v^{p-1}-d(z,s)v^p\bigr)-c v^p=0.
                                                               \tag{H9.119}
\]
Its partial derivative in $z$ at $z=v=0$ is one. Hence
this chart is smooth over the parameter base and is
normal. It has the desired fraction field. There is
exactly one point above each pole: a simple pole cannot
be the pole of $u^p-u$, so the extension has degree $p$;
at a point above the pole, the valuation equation
$p\operatorname{ord}(w)=-e$ forces $e=p$ and
$\operatorname{ord}(w)=-1$. Thus $v$ is a uniformizer,
and (H9.119) is the normalization near the only point
above that pole.

These charts and the \'etale complement prove that the
normalization $\mathscr C\to S$ is a smooth proper
family of connected curves. Its finite map to
$\mathbf P^1\times S$ is flat of degree $p$: both
source and target are regular of equal dimension, so
the finite-morphism flatness criterion applies.
The same argument works with a smooth parameter base
of any dimension while the simple poles remain
disjoint and their residues remain units.

For completeness the ramification and genus can be
computed directly. A nonidentity deck transformation
$w\mapsto w+\alpha$, $\alpha\in\mathbf F_p^\times$,
acts on $v$ by
\[
 v\longmapsto\frac{v}{1+\alpha v},\qquad
 \operatorname{ord}_v\left(\frac{v}{1+\alpha v}-v\right)=2.
\]
The lower ramification groups therefore have order
$p$ in levels $0,1$ and are trivial thereafter. The
different exponent at every pole is $2(p-1)$.
Riemann--Hurwitz gives
\[
 2g(\mathscr C_s)-2=-2p+2(b+2)(p-1),
                      \qquad g=(p-1)(b+1).
\]
The pullback pencil
$\rho:\mathscr X=\widetilde A\times_{\mathbf P^1}\mathscr C\to S$
is smooth and proper: the cover is \'etale over every
old critical value, and over all its branch values
the original pencil is smooth. Its map
$Q:\mathscr X\to\widetilde A\times S$ is finite flat
of degree $p$.

\paragraph{The entire new middle cohomology is a Fourier family.}
Fix $\ell\ne p$ and a nontrivial additive character
$\psi:\mathbf F_p\to\overline{\mathbf Q}_\ell^\times$.
Write $H=H^{n-1}(A,\overline{\mathbf Q}_\ell)$ and let
$V$ be the variable pencil local system of rank $r>0$;
write $q=|\Sigma|$. For a nontrivial deck character
indexed by $\alpha\in\mathbf F_p^\times$, the covering
character sheaf is
\[
             \mathcal L_\psi\bigl(\alpha a(t)+\alpha s/t\bigr).
\]
Set $x=1/t$, and use middle extensions on
$\mathbf A^1_x$ to define
\[
 \mathcal G_\alpha=\mathcal L_\psi(\alpha a(1/x)),\qquad
 \mathcal F_\alpha=(\operatorname{inv}^*V)\otimes\mathcal G_\alpha.
                                                               \tag{H9.120}
\]
The displayed tensor notation means the middle
extension of the tensor product on its common lisse
open set. The old critical points in the $x$ coordinate
are disjoint from the finite poles of $a(1/x)$.

The two outer terms of the middle-degree Leray
filtration are deck invariant, just as in H9-T20.
Its $(1,n-1)$ term is permanent, and condition~A
excludes extra singular-fibre stalks. At the
Artin--Schreier poles the original pencil sheaves
are lisse, whereas the nontrivial character has
Swan conductor one and no inertia invariants. Thus
the exact character part of $R^n\rho_*$ is
\[
 (R^n\rho_*\overline{\mathbf Q}_\ell)_\alpha
  =[s\mapsto\alpha s]^*
     \left(
        H\otimes\operatorname{NFT}_\psi(\mathcal G_\alpha)
           \ \oplus\ \operatorname{NFT}_\psi(\mathcal F_\alpha)
     \right)\bigm|_{\mathbf G_m}.
                                                               \tag{H9.121}
\]
Indeed after $x=1/t$ the parameter factor is exactly
$\mathcal L_\psi(\alpha sx)$. At $x=\infty$, the
function $a(t)$ is regular at $t=0$; for $s\ne0$
the parameter factor has a simple wild pole there,
so its invariants vanish. Therefore parabolic
cohomology equals the compactly supported
$H^1(\mathbf A^1_x,-)$ used in the definition of
$\operatorname{NFT}$. At the old transvection
points the sheaves in (H9.120) already retain the
middle-extension stalks. This proves the relative
Fourier identification with the correct extensions.
The constant ambient term in (H9.121) is retained,
not assumed to vanish.

\paragraph{Both Fourier terms have infinite irreducible monodromy.}
The rank-one sheaf $\mathcal G_\alpha$ has $b+1$
finite simple wild poles, at $x=0$ and $x=1/b_j$.
It is unramified at $x=\infty$. It is an irreducible
Fourier sheaf: its finite wild singularity prevents
it from being $\mathcal L_\psi(cx)$.
The same is true of $\mathcal F_\alpha$ because
tensoring the irreducible $\operatorname{inv}^*V$
with a rank-one sheaf preserves irreducibility.
Katz's Theorem~7.3.8(3) therefore makes both Fourier
transforms irreducible.

Each finite wild pole of $\mathcal G_\alpha$ has
Swan conductor one and drop one. Each such pole
of $\mathcal F_\alpha$ has Swan conductor $r$ and
drop $r$; each old critical point has Swan zero
and drop one. Both input infinity representations
are trivial, of ranks one and $r$ respectively.
Katz's rank formula and the local-invariants
quotient proved in H9-T21 give
\[
 \begin{split}
 \operatorname{rank}\operatorname{NFT}_\psi(\mathcal G_\alpha)
       &=2(b+1),\\
 \operatorname{NFT}_\psi(\mathcal G_\alpha)(0)
       &=J(1,2)\oplus J(1,1)^{\oplus2b},\\
 \operatorname{rank}\operatorname{NFT}_\psi(\mathcal F_\alpha)
       &=q+2r(b+1),\\
 \operatorname{NFT}_\psi(\mathcal F_\alpha)(0)
       &=J(1,2)^{\oplus r}
                      \oplus J(1,1)^{\oplus(q+2rb)}.
 \end{split}                                                  \tag{H9.122}
\]
Each transform is lisse on $\mathbf G_m$ by
Lemma~7.3.9(3), because twisting its input by
$\mathcal L_\psi(zx)$ gives slope one at infinity
when $z\ne0$. The nontrivial size-two blocks in
(H9.122) make the images infinite. The
finite-index-normal-core argument excludes every
finite-orbit vector in each irreducible transform,
and also in a direct sum of its copies. Hence
(H9.121) has no finite-orbit vector for any
nontrivial deck character.

Every algebraic middle class on the geometric
generic fibre is defined over a finite extension
of $k(s)$ and has a finite monodromy orbit.
All its nontrivial deck parts therefore vanish.
For an actual Chow cycle $Z$,
$Z-p^{-1}Q^*Q_*Z$ is globally homologically
trivial. H9-T12 gives admissibility for the entire
homologically trivial generic group, and the
global projection formula proves (H9.118).
Surjectivity of corestriction over $\mathbf Q$
gives its exact radical. These are the same
actual-cycle arguments as H9-T19 and H9-T21;
there is no assertion that the Chow norm kernel
is zero.

\paragraph{The explicit simple-pole parameter stratum.}
One can freely vary the coefficients and poles in
\[
       a_1 t+a_0+\sum_{i=0}^{b}\frac{d_i}{t-e_i},
 \qquad a_1d_0\cdots d_b\ne0,
                                                               \tag{H9.123}
\]
where the $e_i$ are distinct and avoid $\Sigma$.
This is a smooth irreducible parameter open; it
parametrizes the indicated covers with one pole
fixed at infinity and all remaining poles and
nonzero residues free. The local equation
(H9.119) proves smoothness of its normalized
curve family and finite flatness of degree $p$.
The pullback-pencil trace kernel is lisse.
The one-parameter family (H9.117) is obtained
by $a_1=1$, $a_0=0$, $e_0=0$, $d_0=s$ and
the chosen $b_j,c_j$ for the other terms.
Its kernel has no finite-orbit vector, so the
full parameter kernel has none either.
The same conclusion and height identity follow
over the algebraic closure of the function
field of this parameter open. No density claim
about another moduli space is needed here.

\paragraph{Characteristic two: a full double-cover parameter space.}
Now let $p=2$, fix $g\ge1$, and put
$\mathcal L=\mathcal O_{\mathbf P^1}(g+1)$.
Consider the affine space
\[
 \mathscr P_g=H^0(\mathbf P^1,\mathcal L)
                      \oplus H^0(\mathbf P^1,\mathcal L^2)
\]
with parameters $(h,f_2)$ and the universal
finite locally free rank-two algebra
\[
       \mathcal O\oplus\mathcal L^{-1},
                         \qquad y^2+h y=f_2.
                                                               \tag{H9.124}
\]
Its relative spectrum is proper over the
parameter space because it is finite over
$\mathbf P^1\times\mathscr P_g$.
Let $U_g$ be the open where this curve is smooth,
geometrically connected, and unramified above
$\Sigma$. The genus of each such curve is $g$:
\[
 \chi(\mathcal O_C)
    =\chi(\mathcal O)+\chi(\mathcal L^{-1})
    =1-g.
\]
The smooth-curve and unramified conditions are
open by the relative Jacobian and ramification
criteria; connectedness is open in the smooth
proper family. Thus $U_g$, if nonempty, is a
smooth irreducible open of the affine space.

It is nonempty and contains the already proved
Fourier subfamily. Use (H9.117) with $b=g-1$
and let
\[
           H_0(t)=t\prod_{j=1}^{g-1}(t-b_j).
\]
Writing $y=H_0(t)w$ gives
\[
 y^2+H_0(t)y
       =H_0(t)^2\left(t+
           \sum_{j=1}^{g-1}\frac{c_j}{t-b_j}+
                                          \frac{s}{t}\right).
                                                               \tag{H9.125}
\]
The right side is a polynomial of degree at most
$2g+1$. Homogenize $H_0$ to degree $g+1$ and
the right side to degree $2g+2$. These are
sections in (H9.124), depending regularly on
$s\in\mathbf G_m$. At each finite branch
point $h$ has a simple zero and $f_2$ also
has a simple zero with nonzero coefficient;
the same holds at infinity in its local
line-bundle trivialization. The Jacobian
criterion, or (H9.119), proves that these
curves are precisely the smooth normalized
Artin--Schreier curves already considered.
They are unramified over $\Sigma$. Thus this
is an actual morphism $\mathbf G_m\to U_g$.

The universal pullback pencil over $U_g$ is
smooth and proper, with a finite flat
degree-two map to $\widetilde A\times U_g$.
Its middle trace kernel is a lisse direct
summand. Its restriction to (H9.125) has no
finite-orbit vector by (H9.121)--(H9.122).
Hence the entire universal kernel has none.
With constants $\Omega=\overline{k(U_g)}$,
every algebraic kernel class on its geometric
generic fibre is zero. The actual Chow norm
projection gives (H9.118) with $p=2$.

This is the full double-cover parameter space,
not just a claimed dense stratum. Indeed, for
any finite degree-two map $C\to\mathbf P^1$
from a smooth connected genus-$g$ curve,
the map is finite flat and its rank-two
algebra $E=\pi_*\mathcal O_C$ contains the
unit line $\mathcal O$ as a subbundle. Its
quotient is a line bundle $M$. The equality
$\chi(E)=1-g$ forces $\deg M=-(g+1)$.
Thus $M\simeq\mathcal L^{-1}$, and
\[
 \operatorname{Ext}^1(\mathcal L^{-1},\mathcal O)
          =H^1(\mathbf P^1,\mathcal L)=0.
\]
Choosing a splitting gives
$E\simeq\mathcal O\oplus\mathcal L^{-1}$;
its multiplication is exactly a monic
relation $y^2+h y=f_2$ with $h,f_2$ as in
(H9.124). Hence every such cover is represented.
The unramified condition over the nonempty
$\Sigma$ excludes inseparable maps. This
justifies the geometric-generic hyperelliptic
scope without an external density or
Hurwitz-space irreducibility theorem.

\paragraph{The constructed characteristic-two stratum is ordinary.}
In the family (H9.125), the $g$ finite branch
points are $0,b_1,\ldots,b_{g-1}$, with one
more branch point at infinity. The pullbacks
of
\[
       d\log t,\quad d\log(t-b_1),\ldots,
                                      d\log(t-b_{g-1})
\]
are holomorphic on $C$. At any of their base
simple poles the pullback differential has
order $2(-1)+2=0$, since the ramification
index and different exponent are both two;
elsewhere it has no pole. They are independent
because the corresponding rational
differentials on $\mathbf P^1$ have distinct
finite poles and separable pullback is
injective. Since the genus is $g$, they form
a basis of $H^0(C,\Omega_C^1)$. Cartier
fixes each logarithmic differential and
commutes with separable pullback, so it is
bijective on this space. This proves
ordinarity of the explicit simple-pole
family. No assertion about density of this
stratum was used in the preceding full
parameter-space argument.
\end{proof}

\paragraph{Scope and sources.}
The exact Fourier source is the directly read
Katz ESDE material recorded in H9-T21:
Theorem~7.3.8(3), Lemma~7.3.9 and
Corollary~7.4.3.1 with Theorem~7.4.4, or
equivalently Theorem~7.5.4. The wild local
normalization, ramification orders, genus,
both Fourier blocks, and double-cover
parameter argument are proved explicitly
here. The conclusions are for geometric
generic parameters. They do not extend the
sign to every special hyperelliptic cover:
the positive and negative special-cover
directions of H9-T17 remain valid.
